\documentclass{standalone}
\usepackage{tikzplot}

\begin{document}

% begin_snippet
\begin{tikzpicture}
  \axes[grid] (-1:8,-4:4);
  \foreach \i in {-10,...,10}{
    \tikzmath{
      % 定义两圆心, 半径及切点
      \Oa{1} = 2; \Oa{2} = 0; \ra = 2;
      \Ob{1} = 3; \Ob{2} = 0; \rb = 3;
      % 以 A 为反演中心 r 为半径作反演变换
      % 圆 Oa 变换为直线 xa = 0.5 r^2 / ra 
      % 圆 Ob 变换为直线 xb = 0.5 r^2 / rb
      \A{1} = 0; \A{2} = 0; \r = 6;
      \xa = 0.5*(\r)^2/\ra; 
      \xb = 0.5*(\r)^2/\rb;
      % 作与两直线相切的圆
      \d = \xa-\xb;
      \P{1} = \xa; \P{2} = \i*\d; % 直径端点
      \Q{1} = \xb; \Q{2} = \P{2}; % 直径端点
      % 圆上一点
      % 注意: 反演变换后原来的直径不再是直径, 需要通过过三点的圆来构造
      \R{1} = (\P{1}+\Q{1})/2; 
      \R{2} = (\P{1}-\Q{1})/2+\P{2};
    }
    \tikzset{
      plane/inverse={\A,(\r),\P,\Pa},
      plane/inverse={\A,(\r),\Q,\Qa},
      plane/inverse={\A,(\r),\R,\Ra},
      plane/circumcenter={\P,\Q,\R,\C},
      plane/length={\C,\P,\ro},
      plane/circumcenter={\Pa,\Qa,\Ra,\Ca},
      plane/length={\Ca,\Pa,\rc},
    }
    \coordinate (O1) at (\Oa{1},\Oa{2});
    \coordinate (O2) at (\Ob{1},\Ob{2});
    \coordinate (A) at (\A{1},\A{2});
    \coordinate (C) at (\C{1},\C{2});
    \coordinate (P) at (\P{1},\P{2});
    \coordinate (Q) at (\Q{1},\Q{2});
    \coordinate (R) at (\R{1},\R{2});
    \coordinate (C') at (\Ca{1},\Ca{2});
    \coordinate (P') at (\Pa{1},\Pa{2});
    \coordinate (Q') at (\Qa{1},\Qa{2});
    \coordinate (R') at (\Ra{1},\Ra{2});

    \draw[teal] (O1) circle (\ra);
    \draw[violet] (O2) circle (\rb);
    % \draw[purple] (C) circle (\ro);
    \draw[red] (C') circle (\rc);
  }
\end{tikzpicture}
% end_snippet

\end{document}